\section{Background and Related Work} % <1.5 pages

\subsection{GPU Execution model}

% Background on GPU execution model
Basic understanding of the GPU execution model is essential to understand this work. CUDA~\cite{cuda} uses Single Instruction, Multiple Data (SIMD) execution model~\cite{lindholm2008nvidia}. Each GPU kernel launches a global grid of threads to process data concurrently. These threads are grouped into blocks, also known as Cooperative Thread Arrays (CTAs). Each CTA utilizes a localized shared memory to cooperatively compute a specific tile of data. Recent architectures further introduce a abstraction layer called Thread Block Clusters (or Cooperative Group Arrays, CGAs), which groups multiple CTAs together, enabling cooperate via Distributed Shared Memory (DSMEM)~\cite{nvidia2022hopper}. At the hardware level, CTAs are dispatched to Streaming Multiprocessors (SMs), where threads are bundled into groups of 32, known as warps. While individual warps are scheduled and operate independently, threads within the same warp execute the same instruction simultaneously.

% Specialized hardware: tensor core and TMA engine / pipeline kernel
As GPU architectures evolve, core computational workloads are increasingly offloaded to specialized on-chip accelerators. General-purpose SMs are transitioning from performing the primary computation to orchestrating these specialized units. Mature accelerators include Tensor Core for Matrix-Matrix Multiply-Accumulate (MMA), and Tensor Memory Accelerator (TMA). The recently announced integration with Language Processing Unit (LPU) for LLMs also reflects this trend~\cite{LPU}. To efficiently exploit this heterogeneous on-chip setup, warp specialization has emerged as a standard programming paradigm. Specific warps are dedicated to each stages of an execution pipeline, managing asynchronous calls to the hardware accelerators. High-performance GEMM implementations, such as those in the NVIDIA CUTLASS library~\cite{cutlass}, rely on this paradigm to optimize the overall computation.

% Stream, Green context and MiG
% Another aspect of the execution model is the deployment of kernels to GPU resources in the CUDA programming paradigm. In CUDA, kernel launches are dispatched to specific device streams, where multiple streams on the same device can execute concurrently if resource is sufficient. In a multi-GPU setup, streams are tied to specific GPUs, dictating device-level resource selection. Furthermore, even within a single GPU, hardware resources can be explicitly partitioned using mechanisms such as Multi-Instance GPU (MIG)~\cite{mig} or CUDA Green Contexts~\cite{cuda}. These partitioned resources can be bound to specific streams, granting fine-grained control over kernel execution. For simplicity throughout this paper, we refer to each of these segmented resource pools simply as a "GPU" when describing kernel placement.

\subsection{Inter-operator Optimizations}

Large amount of existing work focuses on improving GPU execution efficiency through the co-optimization of consecutive kernels (inter-operator optimization). The boundaries between discrete kernels often introduce performance bottlenecks, including kernel launch overheads~\cite{shen2020nimble}, redundant global memory round-trips~\cite{le2023welder}, and wave quantization effects, where the final wave of scheduled thread blocks fails to fully utilize the available resource. Furthermore, individual kernels have different resource utilization profiles, typically categorized as compute-bound, memory-bound, or communication-bound. By fusing or concurrently executing these operators, it is possible to overlap mismatched resource demands, achieving higher overall hardware utilization.

% Intra-layer: MEGA kernel / kernel fusion on breaking kernel boundary (Control flow)
To directly address inter-operator inefficiencies, kernel fusion approaches this by systematically merging multiple consecutive operations into a single kernel~\cite{chen2018tvm, li2022automatic}. Pushing this optimization paradigm to its extreme results in the mega-kernel approach. In a mega-kernel design, massive computational sub-graphs, ranging from complex attention mechanisms~\cite{dao2022flashattention} to entire multi-GPU model inference pipelines~\cite{wu2024mirage,cheng2025mpk}, are encapsulated within a single GPU kernel launch. However, this extreme form of fusion requires re-compilation of the whole workflow, as well as sophisticated in-kernel runtimes to manage intra-kernel synchronization, warp-level task scheduling, and decentralized hardware resource allocation to prevent SM underutilization~\cite{ma2020rammer}. 

% Micro pipelining in this space 
Besides, inter-operator pipelining offers a complementary approach by overlapping consecutive execution stages. Unlike coarse-grained inter-layer parallelism, these techniques decompose individual operators into fine-grained micro-batches or tiles to enable concurrent processing~\cite{zheng2017versapipe,  cusync2024, zhang2025hytis}. This intra-layer pipelining also serves as a primary mechanism for achieving communication-computation overlap~\cite{zhu2025nanoflow, chang2024flux, wang2023overlap}. 

However, these methods have primarily been deployed as temporal pipelining on single devices to improve overall hardware utilization. The performance benefits of some of those approaches also rely on processing multiple batches of data concurrently, for example, by overlapping the communication of a previous batch with the computation of the current one. The use of fine-grained inter-operator pipelining spatially across multiple devices for reducing single-batch input latency remains largely unexplored. Historically, this is mainly because pipelining introduces pipeline bubbles and requires frequent inter-stage synchronizations, limiting its effectiveness on latency reduction~\cite{zheng2022alpa}.

Very recently, hardware-centric studies such as Kitsune~\cite{davies2025kitsune} have begun exploring the hardware architecture potential to support spatial pipelines. Our work comes from a different angle: we propose a pure software technique designed to exploit the capabilities of modern shared-memory multi-GPU systems, focusing on using fine-grain spatial pipelines to reducing single-batch query latency.
